% Table: Effect of A-level economics on take-up, earnings, and lifetime income (Chapter 4)
\begin{table}[htbp]
\centering
\caption{Effect of A-level economics on take-up, earnings, and lifetime income}
\label{tab:main}
\begin{tabular}{l cccc}
\toprule
                                      & OLS        & Reduced    & IV         & IV           \\
                                      &            & form       & (linear)   & (saturated)  \\
                                      & (1)        & (2)        & (3)        & (4)          \\
\midrule
\multicolumn{5}{c}{\textbf{Panel A: Takes A-level economics (First stage) }} \\
\midrule
A-level economics available           &            & 0.057\sym{***}& 0.057\sym{***}& 0.058\sym{***}  \\
                                      &            & (0.002)       & (0.002)       & (0.002)         \\
\addlinespace
First-stage $F$                       &            &            & 36.3         & 11.1           \\
Lee et al.\ $tF$                      &            &            & 3.7         & 5.8           \\
\midrule
\multicolumn{5}{c}{\textbf{Panel B: Log earnings at age 28}} \\
\midrule
Takes A-level economics               & 0.175\sym{***}&            & 0.135\sym{*}& 0.178\sym{**}  \\
                                      & (0.008)       &            & (0.070)       & (0.072)         \\
\addlinespace
A-level economics available           &            & 0.009\sym{***}&            &              \\
                                      &            & (0.004)       &            &              \\
\midrule
\multicolumn{5}{c}{\textbf{Panel C: Lifetime earnings (PV at age 18, \pounds)}} \\
\midrule
Gross lifetime earnings gain          &            &            & XX\sym{***}& XX\sym{***}  \\
                                      &            &            & (XX)       & (XX)         \\
\addlinespace
Net lifetime earnings gain            &            &            & XX\sym{***}& XX\sym{***}  \\
                                      &            &            & (XX)       & (XX)         \\
\addlinespace
Fiscal return                         &            &            & XX\sym{***}& XX\sym{***}  \\
                                      &            &            & (XX)       & (XX)         \\
\midrule
School FE                             & Yes        & Yes        & Yes        & Yes          \\
Cohort FE                             & Yes        & Yes        & Yes        & Yes          \\
Individual controls                   & Yes        & Yes        & Yes        & Saturated    \\
\addlinespace
Observations                          & 368,257         & 368,257         & 368,262         & 367,937           \\
\bottomrule
\end{tabular}
\begin{minipage}{\textwidth}
\smallskip
\footnotesize
\textit{Notes:} Panel~A reports the first-stage relationship between A-level economics availability and take-up. Panel~B reports the effect of taking A-level economics on log real earnings at age~28 (OLS in column~1; instrumented by availability in columns~3--4); column~2 reports the reduced-form effect of availability on earnings. Panel~C reports IV estimates of the effect of taking A-level economics on lifetime earnings, in present value at age~18. Gross earnings are discounted using UK Treasury Green Book rates (3.5\% for the first 30~years, 3.0\% thereafter). Net earnings deduct income tax, National Insurance, and income-contingent student loan repayments under Plan~2 rules. The fiscal return is the additional tax and NI revenue plus reduced loan write-off cost. Lifetime earnings profiles are constructed following the forward-simulation approach of \citet{britton2020degrees}. Column~4 uses a saturated specification that interacts the instrument with 120 coarse cells defined by sex, ethnicity, FSM status, and GCSE attainment tercile, following \citet{blandhol2022tsls}. Standard errors clustered at the school level in parentheses. Sample: 2002--2009 GCSE cohorts with non-zero earnings at age~28. \sym{*} $p<0.10$, \sym{**} $p<0.05$, \sym{***} $p<0.01$.
\end{minipage}
\end{table}
